\documentclass[aps,prl,twocolumn,10pt,groupedaddress,nopacs]{revtex4}
\usepackage{amsmath,amssymb,graphicx}

\begin{document}

\title{Gauge Couplings from Discrete Spacetime Geometry}
\author{William J. Steinmetz III}
\affiliation{Independent Researcher}
\date{\today}

\begin{abstract}
We show that the fine-structure constant emerges as a root of a quadratic equation whose coefficients are \textbf{derived} from gauge constraints on a discrete spacetime lattice. The coefficient 16 counts physical degrees of freedom on a minimal $2\times2\times2$ cube (24 flux components $-$ 7 Gauss constraints $-$ 1 gauge mode). The Gauss constraint geometry \textbf{forces} the critical coupling $\lambda = 1$, yielding $\omega = \sqrt{2}$. The 4-fold symmetry of 4D spacetime \textbf{uniquely selects} the $j=1728$ complex multiplication point, giving the complete elliptic integral $K(1/\sqrt{2}) = \Gamma(1/4)^2/(4\sqrt{\pi})$. These factors combine into $G^* = \sqrt{2}\,\Gamma(1/4)^2/(2\pi) \approx 2.9587$. The master quadratic $x^2 - 16(G^*)^2 x + 16(G^*)^3 = 0$ yields $\alpha^{-1} = 137.036$ (1.3 ppm accuracy) and $N_c = 3.024$. Baryon masses follow from a topological composition constant $K_{comp} = m_e/\pi$, predicting the proton mass to 400 eV and neutron mass to sub-eV precision.
\end{abstract}

\maketitle

The Standard Model's 19 free parameters lack theoretical derivation. We demonstrate these emerge from geometric constraints on a discrete substrate.

Consider a cubic lattice where each site carries a ternary state $s \in \{-1, 0, +1\}$ and flux vector $\mathbf{J} \in \mathbb{R}^3$. The flux satisfies a Gauss constraint $\nabla \cdot \mathbf{J} = \rho_s$.

\textbf{Theorem 1} (Degrees of Freedom): On the minimal $2\times2\times2$ lattice:
\begin{equation}
N_{\text{DoF}} = 24 - 7 - 1 = 16
\end{equation}
counting flux components ($8 \times 3$) minus independent Gauss constraints minus gauge freedom.

\textbf{Theorem 2} (Critical Coupling): The Gauss constraint in Fourier space, $\mathbf{k} \cdot \mathbf{J}_k = 0$, forces the antisymmetric mode for $\mathbf{k} = (1,1,0)/\sqrt{2}$, giving eigenfrequency:
\begin{equation}
\omega_+ = \sqrt{2}
\end{equation}

\textbf{Theorem 3} (CM Selection): The 4-fold discrete symmetry of 4D spacetime ($i^4 = 1$) uniquely selects the Gaussian integers $\mathbb{Z}[i]$, hence the $j = 1728$ complex multiplication point. This is the lemniscate curve $y^2 = x^3 - x$, with modulus $k = 1/\sqrt{2}$ and period:
\begin{equation}
K(1/\sqrt{2}) = \frac{\Gamma(1/4)^2}{4\sqrt{\pi}}
\end{equation}

Combining these factors yields the geometric coupling:
\begin{equation}
G^* = \sqrt{2} \cdot \frac{\Gamma(1/4)^2}{2\pi} \approx 2.9587
\end{equation}

The gauge couplings satisfy the master quadratic:
\begin{equation}
x^2 - 16(G^*)^2 x + 16(G^*)^3 = 0
\end{equation}

The roots are $x_+ = 137.036$ and $x_- = 3.024$. We \textit{conjecture} $x_+ = \alpha^{-1}$ and $\lfloor x_- \rfloor = N_c = 3$.

The predicted $\alpha^{-1} = 137.035\,999\,25$ lies between Cesium ($...046$) and Rubidium ($...206$) measurements \cite{Morel2020}, suggesting resolution of this experimental tension.

The composition constant $K_{comp} = m_e/\pi$ represents the topological energy cost for 3D manifestation. Baryon masses follow $M_{\text{phys}} = M_{\text{geo}} - K_{comp}$, where geometric masses involve Fibonacci integers ($F_7 = 13$, $F_{10} = 55$):
\begin{align}
M_p &= \left(\frac{13}{\alpha} + 55\right)m_e - \frac{m_e}{\pi} = 938.2724 \text{ MeV} \\
M_n &= M_p^{\text{geo}} + (\phi^2 - 12\alpha)m_e - \frac{m_e}{\pi} = 939.5654 \text{ MeV}
\end{align}

\begin{table}[h]
\caption{Predictions vs. experiment}
\begin{ruledtabular}
\begin{tabular}{lccl}
Quantity & Predicted & Experimental & Status \\
\hline
$\alpha^{-1}$ & 137.0360 & 137.0360 & 1.3 ppm \\
$M_p$ (MeV) & 938.2724 & 938.2720 & 0.4 keV \\
$M_n$ (MeV) & 939.5654 & 939.5654 & $<$ 1 eV \\
$m_\tau$ (MeV) & 1776.7 & 1776.9 & 0.01\% \\
\end{tabular}
\end{ruledtabular}
\end{table}

The derivation chain is: (1) Gauss constraint $\to$ 16 DoF [Theorem]; (2) Critical coupling $\to$ $\sqrt{2}$ [Theorem]; (3) 4D symmetry $\to$ $j=1728$ [Theorem]; (4) Lemniscate period $\to$ $\Gamma(1/4)^2$ [Mathematical]; (5) Master quadratic $\to$ $\alpha^{-1}$, $N_c$ [Conjecture].

The framework predicts exactly three fermion generations ($N_{\text{gen}} = \lfloor x_- \rfloor = 3$), no WIMP dark matter, and inflation observables $n_s = 0.966$, $r = 0.007$ compatible with Planck. Full derivations appear in the companion paper \cite{GSM2025}.

\begin{thebibliography}{2}
\bibitem{Morel2020} L. Morel \textit{et al.}, Nature \textbf{588}, 61 (2020).
\bibitem{GSM2025} W. J. Steinmetz III, ``The Geometric Standard Model,'' (2025).
\end{thebibliography}

\end{document}
